% !TeX program = xelatex
% ======================================================================
% 全国大学生数学建模竞赛（CUMCM）2026 论文模板
%
% 依据本次对话中已确认的 2026 官方要求搭建：
% 1. A4；页边距上下左右至少 2.5 cm；
% 2. 电子论文第一页为“论文标题 + 摘要 + 关键词”，摘要原则上不超过 1 页；
% 3. 不需要英文摘要；
% 4. 不设置目录；
% 5. 正文不超过 30 页，附录页数不限；
% 6. 电子论文不包含承诺书、编号页；
% 7. 摘要、正文、附录不得出现学校、队员姓名、赛区等身份信息；
% 8. 参考文献前设置“AI工具使用声明”；
% 9. 附录中列出支撑材料文件，并给出完整、可运行的主要源程序；
% 10. PDF 与支撑材料压缩包均需控制在赛事要求的大小范围内。
%
% 推荐：
%   latexmk -xelatex main.tex
%
% 推荐项目结构：
%   main.tex
%   cumcmthesis.cls           % 若使用官方/常用 cumcmthesis 模板
%   figures/
%   code/
%   data/
%   support/
%       AI工具使用详情.pdf
%
% 注意：
% - 本模板“章节名称”采用国赛优秀论文常见写法，并非官方强制章节名称。
% - 比赛正式提交前，请再次核对全国组委会及所在赛区/学校最新通知。
% ======================================================================

% ----------------------------------------------------------------------
% 1. 文档类与字体（独立、可稳定嵌入中文字体）
% ----------------------------------------------------------------------
% fontset=none：不依赖操作系统自带的宋体/黑体等字体；
% 明确使用 Noto CJK 字体，避免 PDF 中“中文消失，只剩数字/斜杠”的问题。
\documentclass[UTF8,fontset=none,a4paper,12pt]{ctexart}

\usepackage[
  a4paper,
  left=2.5cm,
  right=2.5cm,
  top=2.5cm,
  bottom=2.5cm
]{geometry}
\usepackage{amsmath,amssymb,bm}
\usepackage{graphicx}
\usepackage{booktabs}
\usepackage{longtable}
\usepackage{multirow}
\usepackage{listings}
\usepackage{xcolor}
\usepackage{subcaption}
\usepackage{float}
\usepackage[hidelinks]{hyperref}

% ---------- 字体体系：正文宋体风格 + 标题黑体风格 + 英文 Times 风格 ----------
% 中文：
%   正文  -> Noto Serif CJK SC（宋体/明朝体风格）
%   标题  -> Noto Sans CJK SC（黑体风格）
%   代码  -> Noto Sans Mono CJK SC
% 西文：
%   正文  -> Liberation Serif（Times New Roman 兼容风格）
%   标题  -> Liberation Sans
%   代码  -> Liberation Mono
%
% 所有字体均由 XeLaTeX 嵌入 PDF，避免浏览器预览时中文消失。
\setmainfont{Liberation Serif}
\setsansfont{Liberation Sans}
\setmonofont{Liberation Mono}

\setCJKmainfont{Noto Serif CJK SC}
\setCJKsansfont{Noto Sans CJK SC}
\setCJKmonofont{Noto Sans Mono CJK SC}

\setlength{\parindent}{2em}
\linespread{1.18}

% 页码从摘要页开始，页脚居中
\pagestyle{plain}
\pagenumbering{arabic}
\setcounter{page}{1}

% “一、问题重述”“二、问题分析”……
\ctexset{
  section = {
    name = {,、},
    number = \chinese{section},
    format = \centering\sffamily\bfseries\zihao{4},
    aftername = \hspace{0.4em},
    beforeskip = 1.2ex plus .2ex,
    afterskip = 1.2ex plus .2ex
  },
  subsection = {
    format = \sffamily\bfseries\zihao{-4},
    beforeskip = 1.0ex plus .2ex,
    afterskip = .7ex plus .2ex
  },
  subsubsection = {
    format = \sffamily\bfseries\zihao{-4},
    beforeskip = .8ex plus .2ex,
    afterskip = .5ex plus .2ex
  }
}

\graphicspath{{figures/}}
\renewcommand{\arraystretch}{1.15}
\setlength{\tabcolsep}{6pt}

% ----------------------------------------------------------------------
% 3. 代码样式
% ----------------------------------------------------------------------
\lstset{
  basicstyle=\ttfamily\small,
  numbers=left,
  numberstyle=\scriptsize,
  numbersep=8pt,
  frame=single,
  breaklines=true,
  columns=fullflexible,
  showstringspaces=false,
  tabsize=4
}

% ----------------------------------------------------------------------
% 4. 论文标题（电子版匿名）
% ----------------------------------------------------------------------
\title{\sffamily\bfseries\zihao{2} XXXX问题的数学模型}
\author{}
\date{}
\providecommand{\keywords}[1]{\par\vspace{0.5em}\noindent\textbf{关键词：}#1}

% ----------------------------------------------------------------------
% 5. 常用数学命令
% ----------------------------------------------------------------------
\newcommand{\R}{\mathbb{R}}
\newcommand{\E}{\mathbb{E}}
\newcommand{\Var}{\operatorname{Var}}
\newcommand{\argmin}{\operatorname*{arg\,min}}
\newcommand{\argmax}{\operatorname*{arg\,max}}

\begin{document}

% ----------------------------------------------------------------------
% 版式说明（比赛写作时无需改）
% ----------------------------------------------------------------------
% 本模板采用：
% - 中文正文：宋体风格；
% - 中文标题：黑体风格；
% - 英文/数字：Times New Roman 兼容风格；
% - 数学公式：LaTeX 标准数学字体。
%
% 全国组委会对“字号、字体、行距、颜色”没有统一强制规定；
% 因此这里采用国内科技论文常见、阅读性较好的统一字体体系。
% ----------------------------------------------------------------------

% ======================================================================
% 第一页：论文标题 + 摘要 + 关键词
% 官方要求电子论文第一页直接进入摘要页，不含承诺书/编号页。
% 摘要原则上控制在 1 页以内；不写英文摘要。
% ======================================================================
\maketitle

\begin{abstract}
% ---------------- 摘要写作建议 ----------------
% 建议直接回答四件事：
% 1. 研究了什么；
% 2. 各小问分别建立了什么模型、用了什么方法；
% 3. 得到了哪些关键数值/结论；
% 4. 如何检验模型、模型价值是什么。
%
% 避免写成“本文首先……其次……最后……”的目录式流水账。
% 尽量给出关键指标、误差、最优参数、预测结果等定量信息。
% ----------------------------------------------

针对题目所研究的\textbf{核心对象与目标}，本文从\textbf{主要矛盾/关键机制}出发，
建立了由\textbf{模型一}、\textbf{模型二}和\textbf{综合优化模型}组成的统一分析框架。

针对问题一，首先对原始数据进行\textbf{清洗、重构与特征提取}，在若干合理假设下建立
\textbf{模型名称}。通过\textbf{求解算法}得到……，关键结果为……。

针对问题二，在问题一模型基础上引入\textbf{新增因素/约束}，建立\textbf{模型名称}，
并利用\textbf{算法或数值方法}求解。结果表明……。

针对问题三，进一步考虑\textbf{泛化场景/不确定性/优化目标}，构建\textbf{综合模型}。
通过\textbf{敏感性分析、误差分析或稳健性检验}验证模型可靠性，最终得到……。

综上，本文模型具有\textbf{可解释性、稳定性与一定推广能力}，可为……提供定量依据。

\keywords{关键词1\quad 关键词2\quad 关键词3\quad 关键词4}
\end{abstract}

% 官方要求：不设置目录
% \tableofcontents

% ======================================================================
\section{问题重述}
% ======================================================================

\subsection{问题背景}

随着……的发展，……问题逐渐受到关注。题目给出了……数据（或条件），
要求从……角度研究……之间的关系，并进一步给出……方案。
因此，本文需要在充分利用题目信息的基础上，建立具有较强解释性与可计算性的数学模型。

% 写作提示：
% - 只保留与建模有关的背景；
% - 不整段照抄题面；
% - 通常 1--2 段即可。

\subsection{问题提出}

根据题意，本文需要解决以下问题：

\textbf{问题一：} 根据题目给定的……，建立数学模型计算/预测/评价……，并给出……结果。

\textbf{问题二：} 在问题一的基础上，进一步考虑……因素，分析……并求解……。

\textbf{问题三：} 综合考虑……与……，建立优化/评价/预测模型，得到……方案。

\textbf{问题四：} 将模型推广到……场景，讨论模型适用性并提出相应建议。

% 若题目只有三问，删除问题四及后续对应章节。

% ======================================================================
\section{问题分析}
% ======================================================================

\subsection{问题一分析}

问题一的核心是由已知的……确定未知的……。由于……具有……特征，
因此可先对……进行预处理，再利用……关系构建模型。
具体而言，首先……；随后……；最终通过……得到所需结果。

\subsection{问题二分析}

问题二在问题一基础上增加了……约束/变量，使得原问题由……
转化为……问题。为此，需要在前述模型中引入……，
并采用……方法对参数或决策变量进行求解。

\subsection{问题三分析}

问题三需要同时兼顾……、……与……，具有明显的多目标/非线性/不确定性特征。
因此本文拟构建……指标体系，并采用……方法统一衡量各因素，
最后通过……得到综合最优结果。

\subsection{问题四分析}

问题四重点考察模型的泛化能力。本文将在原模型的基础上调整……参数，
分析模型在不同……条件下的变化，并从……角度提出可操作建议。

\subsection{建模思路}

本文总体建模思路如下：
\begin{enumerate}
  \item 对题目数据进行清洗、单位统一、异常值检测与必要的特征提取；
  \item 针对问题一建立基础模型，确定关键变量与参数；
  \item 在基础模型上逐步加入问题二、问题三中的新因素和约束；
  \item 通过数值算法完成模型求解，并利用图表展示结果；
  \item 采用误差分析、敏感性分析或稳健性检验评价模型；
  \item 根据结果给出结论，并讨论模型的推广价值。
\end{enumerate}

% 将下方占位框替换为流程图：
\begin{figure}[H]
  \centering
  \fbox{\rule{0pt}{5.0cm}\rule{0.86\textwidth}{0pt}}
  \caption{总体建模流程图}
  \label{fig:workflow}
\end{figure}

% ======================================================================
\section{模型假设}
% ======================================================================

为了突出问题的主要矛盾并保证模型可求解，作如下假设：

\begin{enumerate}
  \item \textbf{假设一：} 题目所给数据真实、可靠，能够反映研究对象的基本规律；
  \item \textbf{假设二：} 在研究时间或空间范围内，未明确给出的外部因素对结果影响较小；
  \item \textbf{假设三：} 模型涉及的主要参数在相应研究区间内保持稳定；
  \item \textbf{假设四：} 对个别缺失或异常数据进行合理处理后，不改变整体变化趋势。
\end{enumerate}

% 正式比赛时必须改成与具体题目强相关的假设。
% 不建议为了凑数量写“空气阻力忽略”等与实际模型无关的套话。

% ======================================================================
\section{符号说明}
% ======================================================================

本文主要符号及其含义如表~\ref{tab:symbols} 所示。

\begin{table}[H]
  \centering
  \caption{主要符号说明}
  \label{tab:symbols}
  \begin{tabular}{ccc}
    \toprule
    符号 & 含义 & 单位 \\
    \midrule
    $x_i$ & 第 $i$ 个样本或决策变量 & -- \\
    $y_i$ & 第 $i$ 个观测值 & -- \\
    $\hat y_i$ & 第 $i$ 个预测值 & -- \\
    $\bm{\theta}$ & 模型参数向量 & -- \\
    $J$ & 目标函数值 & -- \\
    $w_i$ & 第 $i$ 个指标的权重 & -- \\
    \bottomrule
  \end{tabular}
\end{table}

% 建议只列高频、重要、容易混淆的符号。
% 简单且只出现一次的变量，可直接在公式后解释。

% ======================================================================
\section{数据预处理}
% ======================================================================

\subsection{数据清洗}

首先对题目所给数据进行完整性与合理性检查，主要包括：
\begin{enumerate}
  \item 检查缺失值、重复值及明显异常值；
  \item 统一不同数据表中的单位、时间尺度与空间尺度；
  \item 根据数据特征采用插值、平滑或异常值修正方法；
  \item 对后续建模所需变量进行重构与特征提取。
\end{enumerate}

\subsection{描述性统计}

对关键变量计算均值、标准差、极值及相关性，并通过可视化识别主要变化规律。

\begin{figure}[H]
  \centering
  \fbox{\rule{0pt}{4.5cm}\rule{0.78\textwidth}{0pt}}
  \caption{关键变量描述性分析示意图}
  \label{fig:eda}
\end{figure}

% 若题目本身不存在数据预处理，可删掉整个第五章，
% 后续章节编号会自动调整。

% ======================================================================
\section{模型的建立与求解}
% ======================================================================

% ----------------------------------------------------------------------
\subsection{问题一模型的建立与求解}

\subsubsection{模型建立}

设问题一的决策变量为
\[
\bm{x}=(x_1,x_2,\ldots,x_n)^\mathsf{T}.
\]

根据题目中的……关系，建立目标函数
\begin{equation}
  \min_{\bm{x}} J(\bm{x})
  = \sum_{i=1}^{m} w_i L_i(\bm{x})
  + \lambda R(\bm{x}),
  \label{eq:q1-objective}
\end{equation}
其中，$L_i(\bm{x})$ 表示……，$w_i$ 表示……，$R(\bm{x})$ 表示……。

相应约束条件为
\begin{equation}
\begin{cases}
  g_j(\bm{x})\le 0, & j=1,2,\ldots,p,\\
  h_k(\bm{x})=0, & k=1,2,\ldots,q.
\end{cases}
\label{eq:q1-constraint}
\end{equation}

\subsubsection{参数确定}

说明参数来自：
题目直接给定、数据拟合、文献取值、经验设定，或通过反演/优化求得。
若参数需要拟合，应说明损失函数、拟合区间与最终参数值。

\subsubsection{模型求解}

根据模型特点，采用……算法进行求解。主要步骤为：
\begin{enumerate}
  \item 输入数据及初始化模型参数；
  \item 计算目标函数或状态变量；
  \item 检查约束条件并更新决策变量；
  \item 判断是否满足终止条件；
  \item 输出最终解及相关评价指标。
\end{enumerate}

\subsubsection{结果分析}

问题一的主要结果如表~\ref{tab:q1-result} 所示。

\begin{table}[H]
  \centering
  \caption{问题一主要计算结果}
  \label{tab:q1-result}
  \begin{tabular}{cccc}
    \toprule
    指标 & 结果1 & 结果2 & 说明 \\
    \midrule
    指标1 & -- & -- & -- \\
    指标2 & -- & -- & -- \\
    指标3 & -- & -- & -- \\
    \bottomrule
  \end{tabular}
\end{table}

由表~\ref{tab:q1-result} 可知，……。这表明……，并与实际规律保持一致。

% ----------------------------------------------------------------------
\subsection{问题二模型的建立与求解}

\subsubsection{模型建立}

问题二在问题一基础上进一步考虑……。
因此，引入变量或参数……，将原模型扩展为
\begin{equation}
  \bm{x}^{*}
  = \argmin_{\bm{x}\in\Omega}F(\bm{x};\bm{\theta}),
  \label{eq:q2-model}
\end{equation}
其中，$\Omega$ 为可行域，$\bm{\theta}$ 为模型参数。

\subsubsection{模型求解}

采用……方法求解，并说明搜索区间、步长、初值、迭代次数、终止条件等关键设置。

\subsubsection{结果分析}

给出图表与关键数值，并重点回答：
\textbf{“结果是多少、为什么会这样、是否符合实际规律”。}

% ----------------------------------------------------------------------
\subsection{问题三模型的建立与求解}

\subsubsection{模型建立}

根据问题三要求，在已有模型上进一步引入……因素。
若问题涉及多目标，可构造综合目标函数
\begin{equation}
  S_i=\sum_{j=1}^{K}w_jz_{ij},
  \qquad
  \sum_{j=1}^{K}w_j=1,\quad w_j\ge 0.
  \label{eq:q3-score}
\end{equation}

若不同指标量纲不同，应在加权前进行合理的无量纲化处理。

\subsubsection{模型求解与结果}

根据综合得分或优化结果得到……，并给出排序、最优参数或最终策略。

% ----------------------------------------------------------------------
\subsection{问题四模型的建立与求解}

% 若题目只有三问，删除本节。

\subsubsection{模型推广}

将前述模型推广到……条件下。
通过修改……参数、……边界条件和……约束，即可得到新场景下的模型。

\subsubsection{推广结果与建议}

结合推广结果，从……、……和……三个方面提出建议。

% ======================================================================
\section{模型检验与敏感性分析}
% ======================================================================

\subsection{误差分析}

若模型包含预测或拟合部分，可采用均方根误差
\begin{equation}
  \mathrm{RMSE}
  = \sqrt{
      \frac{1}{n}
      \sum_{i=1}^{n}
      (y_i-\hat y_i)^2
    },
  \label{eq:rmse}
\end{equation}
并结合 MAE、MAPE、$R^2$ 或残差图评价模型精度。

% 请根据题型选择合适指标，不能机械套用 RMSE。

\subsection{敏感性分析}

选择关键参数 $\theta$，在基准值 $\theta_0$ 附近进行扰动：
\begin{equation}
  \theta=(1+\delta)\theta_0,
  \qquad
  \delta\in[-\varepsilon,\varepsilon].
\end{equation}

观察核心输出与最优方案是否发生显著变化，从而判断模型对参数变化的敏感程度。

\subsection{稳健性检验}

可根据题型采用：
随机噪声扰动、Bootstrap 重采样、Monte Carlo 模拟、
不同初值重复求解或不同数据子集交叉验证等方法。

% ======================================================================
\section{模型评价与改进}
% ======================================================================

\subsection{模型优点}

\begin{enumerate}
  \item \textbf{针对性强：} 模型围绕题目核心变量建立，能够直接回答各小问；
  \item \textbf{层次清晰：} 后续模型在前一问基础上逐步扩展，逻辑连贯；
  \item \textbf{可解释性较好：} 主要参数具有明确的实际意义；
  \item \textbf{稳健性较好：} 在合理扰动下主要结论保持稳定。
\end{enumerate}

\subsection{模型不足}

\begin{enumerate}
  \item 为保证模型可求解，对部分现实因素进行了简化；
  \item 部分参数依赖有限数据估计，可能存在一定误差；
  \item 在更复杂场景或更大数据规模下，求解成本可能增加。
\end{enumerate}

\subsection{模型改进}

可从数据精度、模型机制、参数估计、优化算法和场景推广等方面进一步改进。

% ======================================================================
\section{结论}
% ======================================================================

本文围绕题目提出的若干问题建立了相应数学模型，主要结论如下：

\begin{enumerate}
  \item \textbf{问题一：} ……，最终得到……；
  \item \textbf{问题二：} ……，最终得到……；
  \item \textbf{问题三：} ……，最终得到……；
  \item \textbf{问题四：} ……，最终得到……。
\end{enumerate}

综合来看，所建立模型能够较好描述……规律，并可为……提供定量参考。

% ======================================================================
% AI 工具使用声明
% 官方要求：放在“参考文献之前”。
% 二选一保留；若实际使用 AI，请保留第一种，并在支撑材料中提交
% “AI工具使用详情.pdf”。
% ======================================================================

\section*{AI工具使用声明}
\addcontentsline{toc}{section}{AI工具使用声明}

% ---------- 使用过 AI 时保留以下版本 ----------
本参赛队在竞赛过程中使用了AI工具，主要用于
\textbf{语言表达辅助、程序调试、资料整理或其他实际使用环节}，
所有涉及建模、计算、结果与结论的内容均由参赛队成员进行人工审核与核验，
详细使用情况见支撑材料中的《AI工具使用详情.pdf》。

% ---------- 完全未使用 AI 时，删除上段并取消下一行注释 ----------
% 本参赛队在竞赛过程中未使用任何AI工具。

% ======================================================================
% 参考文献
% ======================================================================

\begin{thebibliography}{99}

\bibitem{ref1}
作者.
\newblock 文献题目[J].
\newblock 期刊名, 年份, 卷(期): 页码.

\bibitem{ref2}
作者.
\newblock 书名[M].
\newblock 出版地: 出版社, 年份.

\bibitem{ref3}
机构名称.
\newblock 数据集或网页名称[EB/OL].
\newblock 访问日期：2026-09-XX.

\end{thebibliography}

% ======================================================================
% 附录
% 官方要求：附录页数不限。
% 应包含支撑材料文件列表，以及建模过程中使用的完整、可运行源程序等。
% ======================================================================

\clearpage
\appendix

\section{支撑材料文件列表}

% 请与最终提交的 ZIP/RAR 实际内容保持完全一致。
% 官方题目直接提供的原始附件通常不必重复放入支撑材料。
\begin{longtable}{p{0.34\textwidth}p{0.58\textwidth}}
  \caption{支撑材料文件列表}\label{tab:support-files}\\
  \toprule
  文件名 & 内容说明 \\
  \midrule
  \endfirsthead

  \toprule
  文件名 & 内容说明 \\
  \midrule
  \endhead

  code/q1.py & 问题一完整求解程序 \\
  code/q2.py & 问题二完整求解程序 \\
  code/q3.py & 问题三完整求解程序 \\
  data/processed.csv & 建模过程中生成的处理后数据 \\
  figures/ & 论文中使用的主要结果图 \\
  AI工具使用详情.pdf & AI 工具使用情况、典型交互、采纳与人工核验说明 \\
  \bottomrule
\end{longtable}

% 若未使用 AI，应删除“AI工具使用详情.pdf”这一行。

\section{完整源程序}

% 官方要求的是“完整、可运行”程序。
% 最终比赛版本不要只贴伪代码或删减后的核心片段。
%
% 程序较长时，可按问题拆成多个附录：
%   附录B 问题一程序
%   附录C 问题二程序
%   ……
%
% 下面仅作为格式示例。

\begin{lstlisting}[language=Python,caption={问题一完整程序示例},label={code:q1}]
import numpy as np

def objective(x):
    return np.sum(np.asarray(x) ** 2)

def main():
    x = np.zeros(5)
    result = objective(x)
    print(result)

if __name__ == "__main__":
    main()
\end{lstlisting}

\section{补充计算结果}

% 可放：
% - 大尺寸参数表
% - 大量中间结果
% - 敏感性分析补充图
% - 详细推导
% - 正文因篇幅未展示但用于支撑结论的数据

\begin{table}[H]
  \centering
  \caption{补充参数表}
  \begin{tabular}{ccc}
    \toprule
    参数 & 取值 & 说明 \\
    \midrule
    $\theta_1$ & -- & -- \\
    $\theta_2$ & -- & -- \\
    \bottomrule
  \end{tabular}
\end{table}

% ======================================================================
% 提交前最终检查（仅注释，不会出现在论文中）
%
% [ ] 第一页是否为“标题 + 摘要 + 关键词”
% [ ] 摘要是否原则上控制在 1 页以内
% [ ] 是否删除英文摘要
% [ ] 是否没有目录
% [ ] 正文是否不超过 30 页
% [ ] 页边距是否均不少于 2.5 cm
% [ ] 是否从摘要页开始显示页码
% [ ] 是否完全删除学校、姓名、赛区等身份信息
% [ ] 是否没有承诺书和编号页
% [ ] 所有外部资料是否规范引用
% [ ] AI工具使用声明是否位于参考文献之前
% [ ] 若使用 AI，支撑材料中是否有“AI工具使用详情.pdf”
% [ ] 附录是否列出支撑材料文件
% [ ] 源程序是否完整、可运行，且输出与论文一致
% [ ] 支撑材料文件列表与 ZIP/RAR 实际内容是否一致
% [ ] PDF 文件大小是否符合最新要求
% [ ] 支撑材料压缩包大小是否符合最新要求
% [ ] 是否再次核对全国组委会及所在赛区最新通知
% ======================================================================

\end{document}
